\chapter{Why a Game, Why \riscv, Why Now}
\label{ch:intro}

\section{The problem this book solves}

A first-year student of computing is, in a sense, asked to believe two stories at
once. The first story is told in the computer architecture course: a processor is
a finite collection of registers and a small repertoire of instructions, and
everything a computer does is, at bottom, a long sequence of those instructions
moving numbers between registers and memory. The second story is told in the
programming course: software is built from variables, functions, loops, and the
clean abstractions of a high-level language, and the machine underneath is a
detail best left to compilers. Both stories are true. The difficulty is that a
beginner is rarely given a single object in which both stories are visibly the
same object seen from two sides.

This book provides that object: a small retro video game running on a real
\riscv processor. In the assembly track you will write the game one instruction
at a time and watch a rectangle move because you set a bit. In the C track you
will write the same game as ordinary functions and watch the same rectangle move.
Because both tracks compile to native \riscv machine code and run on the same
platform, the gap between the two stories closes. The compiler stops being a
black box and becomes a translator whose output you can read.

\begin{prgconcept}
The central pedagogical claim of this book is that a visible, audible, native
program collapses the distance between an instruction and its consequence. A
grade is an opinion; a paddle that leaves the screen is a fact. We will use facts.
\end{prgconcept}

\section{The platform in one paragraph}

The platform is \prg, an open educational runtime for \riscv assembly and C games
\citep{prg32repo}. It is important to understand what \prg is \emph{not}: it is
not a CPU instruction emulator that interprets your code on top of another
machine. Your code runs natively, as 32-bit \riscv machine code, either on a
physical ESP32\nobreakdash-C6 board or on Espressif's QEMU firmware target. Games are
distributed as small \emph{cartridges}, files with the extension \fname{.prg32},
which the resident firmware loads and runs through three functions you will write
yourself: an \code{init} called once, and an \code{update} and a \code{draw}
called once per frame. That three-function shape is the entire contract between
your game and the machine, and it is identical whether you write in assembly or
in C. We will return to it in Chapter~\ref{ch:platform}; for now it is enough to
know that the platform is deliberately thin, so that the processor it runs on is
never far out of view.

\section{\riscv as strategic European intellectual property}

It would be possible to teach this material on any number of processors. We have
chosen \riscv deliberately, and the reason reaches well beyond the classroom.

For most of the history of modern computing, the instruction set architecture---
the fundamental contract that defines what a processor's instructions mean---has
been proprietary. The dominant architectures of the desktop, server, and mobile
eras were owned by a small number of companies, and designing a competitive
processor meant either licensing one of those architectures under restrictive
terms or being shut out entirely. An instruction set is the deepest layer of a
computing stack at which a nation or a region can exercise genuine technological
independence, and for decades that layer was not open.

\riscv changed the terms of the question. It is an open standard instruction set
architecture, built on established reduced-instruction-set principles and
governed as a shared specification rather than a private product
\citep{riscvisa}. Anyone may design, manufacture, and sell a \riscv processor
without paying a licensing fee for the architecture itself. That openness is not
merely a convenience for hobbyists; it has become a matter of explicit industrial
strategy, and nowhere more visibly than in Europe.

\begin{prgaside}
The name ``\riscv''---pronounced ``risk-five''---marks the fifth generation of
reduced-instruction-set computing research at the University of California,
Berkeley, where the architecture began as an academic project. That it grew from
a teaching and research instruction set into a strategic industrial standard is a
useful reminder that the line between a classroom exercise and a national
technology base is thinner than it looks.
\end{prgaside}

The European Union has come to treat \riscv as a strategic lever for
technological sovereignty and for greater competition in a global semiconductor
market worth on the order of seven hundred billion United States dollars
\citep{hipeac2026uap}. Europe currently accounts for roughly a tenth of that
market, and the 2023 European Chips Act sets the explicit goal of doubling that
share to a fifth by the end of the decade \citep{hipeac2026uap}. Achieving such a
goal on borrowed, proprietary architectures would leave the continent's most
strategic hardware dependent on decisions made elsewhere; building it on an open
architecture does not.

This is not aspiration on paper alone. Through the EuroHPC Joint Undertaking,
Europe has launched the DARE project---``Digital Autonomy with \riscv in
Europe''---a multi-year, multi-partner effort, backed by funding on the order of
two hundred and forty million euros, to design European processors and
accelerators for high-performance computing and artificial intelligence on a
\riscv foundation \citep{eurohpcdare,daremedium}. DARE builds on a lineage of
earlier European programmes, including the European Processor Initiative, and is
paralleled by industrial efforts such as the TRISTAN project under the Chips
Joint Undertaking, whose partners have assembled Europe's first comprehensive
collection of verified, industry-ready \riscv intellectual property
\citep{hipeac2026uap}. The recurring word in all of these announcements is
\emph{sovereignty}: the capacity to build strategic computing technology without
asking permission.

\begin{prgconcept}
\riscv is open intellectual property: a freely usable instruction set rather than
a licensed product. Europe has made it a centrepiece of its strategy for
technological sovereignty, funding processor and high-performance-computing
programmes built on it. Learning to read and write \riscv is therefore not only a
foundational computing skill but a small act of participation in that strategy.
\end{prgconcept}

For a first-year student the lesson is concrete and encouraging. The instructions
you will trace by hand in Chapter~\ref{ch:asm1}---\texttt{addi}, \texttt{lw},
\texttt{sw}, \texttt{beqz}---are not a teaching dialect invented for a course.
They are the same instructions that European industry and academia are using to
build the processors that the continent intends to depend on. The skill transfers
directly, and it transfers to a technology whose importance is still increasing.

\section{Why retro games, and what \emph{Ready Player One} understood}

If \riscv answers the question ``which processor,'' retro game development
answers the question ``which kind of program.'' The choice is, again, deliberate.

Retro games occupy a pedagogical sweet spot. They are small enough that a
beginner can hold an entire program in their head, yet rich enough to exercise
every concept a first course needs: state stored in memory, input read from the
world, arithmetic that updates that state, branches that make decisions, and a
loop that ties it all together at a fixed frame rate. A game written for a
320-by-200 retro viewport, with a handful of moving rectangles, is simultaneously
a complete and satisfying artefact and a transparent specimen for study. Crucially,
a game gives \emph{immediate, honest feedback}. A sorting routine that is subtly
wrong may still print plausible numbers; a game that is subtly wrong puts the ball
through the paddle, and the error announces itself.

There is also a cultural argument, and it is here that Ernest Cline's novel
\emph{Ready Player One} is more than a passing reference \citep{cline2011}. The
novel imagines a future in which fluency with the games, machines, and idioms of
an earlier computing era becomes the most valuable literacy of all---the key, in
the story quite literally, to a vast inheritance. Stripped of its science-fiction
framing, the book captures something true about computing as a discipline: the
foundational layers do not become obsolete so much as they become invisible, and
the people who can still see them retain a kind of power. The student who
understands how a sprite is drawn pixel by pixel, how a collision is detected with
a few comparisons, and how a frame is timed, understands the substrate on which
all the comfortable modern abstractions are built. Retro game development is, in
this sense, a way of preserving and transmitting exactly the literacy that
\emph{Ready Player One} treats as priceless.

\begin{prgaside}
The aesthetic of this book---the 320-by-200 viewport, the chunky rectangles, the
single-channel beep, the cartridge that you ``insert'' by uploading a
\fname{.prg32} file---is a faithful homage to the home computers and consoles of
the late twentieth century. The constraints are not nostalgia for its own sake.
A tight memory budget and a simple display force a beginner to understand what is
actually happening, in a way that a modern engine with gigabytes of memory and a
forgiving runtime never would.
\end{prgaside}

\section{Both targets, no compromise}

A practical commitment shapes every example in this book: it must run on both the
desktop emulator and the physical board. The emulator, Espressif's QEMU build,
lets you compile, run, and debug graphics on a laptop with no hardware at all; it
is ideal for checking layout, movement, timing, and for stepping through code
with a debugger. The physical ESP32\nobreakdash-C6 board is where the program meets the
real world: real buttons, a real display, a real buzzer, and wireless cartridge
upload \citep{prg32qemu}. The QEMU target is configured as an ESP32\nobreakdash-C3 emulator
profile because that is the \riscv graphics path Espressif maintains, while the
hardware is the ESP32\nobreakdash-C6; the \riscv calling convention and the \prg interface
do not change between them, so the same source builds for both.

\begin{prgtargets}
\textbf{Both targets, every example.} Throughout the book this badge marks code
that has been designed to run unchanged on QEMU and on the ESP32\nobreakdash-C6. Where a
feature is hardware-only (physical buttons, the buzzer, Wi-Fi upload) or
emulator-only (host-side debugging conveniences), the text says so explicitly.
\end{prgtargets}

\section{How the rest of the book is arranged}

Part~I, which this chapter opens, sets the foundation. Chapter~\ref{ch:platform}
introduces game design principles and the \prg framework in detail; it is the one
chapter that every reader, on either track, should read before going further.

Part~II is the computer architecture track, in \riscv assembly. It begins with
your first program, builds through state and input, develops a complete graphical
game, and ends with a capstone you assemble yourself.

Part~III is the computer programming track, in C. It is structured to parallel
Part~II without depending on it: the same three-function shape, the same runtime,
the same two targets, expressed in a high-level language.

Part~IV is a short reunion. Having written games in both languages, you will look
at the two side by side and see, concretely, what a compiler does.

The appendices are the reference shelf: the complete runtime interface,
the performance-measurement workflow, full annotated example programs, a guide to
using GitHub for your coursework, instructions for building the physical board,
and step-by-step environment setup for Windows, Linux, and macOS. The chapters
send you to the appendices rather than interrupting the narrative with reference
tables.

\begin{prgcheck}
Before continuing, you should be able to explain, in a sentence or two each: why
this book uses a game as its teaching object; what it means to say \riscv is an
open instruction set and why Europe treats it as strategic; and what it means
that every example runs on both QEMU and the ESP32\nobreakdash-C6.
\end{prgcheck}
